\section{Answers to the questions as posed}\label{sec:answers}

\paragraph{Can a malign prior hypothesis be stated precisely in a toy?} Yes: \cref{def:mph}. Fix the machine,
the bounds, the block length, the test that marks the handover, and the world. Then ``most of the posterior
is off the world's own continuation'', ``at least a share $\theta$ of it is on continuations in a given set'',
and ``at least a share $\theta$ of it is on continuations that recreate the handover with themselves in the
world's place'' are statements about finitely many rational numbers. They are necessary for calling the prior
malign at that world and not sufficient, since an unpredictable world satisfies them too.

\paragraph{Worlds as programs that call the predictor and run the result: is there a decent amount of quining
out by hypotheses that look for such opportunities and predict well up to them?} Two questions are in this
one (\cref{sec:question}), and the toy answers them differently.

\emph{How far does one such hypothesis reach?} As far as its own prior weight. Without resource bounds, a
program that enumerates the programs making a critical use, spreads its mass over them in proportion to their
prior weight, prints their output up to the use and then a continuation of its choice has, at every use, share
more than $2^{-a-1}\kappa(x)/H$. Here $a$ is its length including the bits that name
$j=\lfloor\log_2(1/H)\rfloor+1$, $H$ is the prior probability of a critical use, and $\kappa(x)$ is the share
of the posterior at the context held by the worlds that make the use (\cref{thm:critical}(a)). It is the
source of the continuation that is run in worlds of total prior probability more than $2^{-a-1}$ times the
average of $\kappa$ over uses. Rarity of critical uses does not reduce this. If the continuation is to be the
program's own code, the same holds with a longer program and for the worlds whose programs fit into a block
(\cref{thm:critical}(b), \cref{thm:search}(c)); without any search, a wrapper gives $2^{-c}\rho^\circ(x)$ at
every use that a fixed test recognises (\cref{thm:quine}). A seeker can also concentrate on a part of the uses
and hold the majority there, at uses of total prior mass up to about its own weight times the fraction that
the genuine users are of the other programs printing those contexts.

\emph{Do most uses at which this is possible end that way?} That does not follow, and it is the question the
argument is about. The share of critical uses that end in a takeover is the average posterior share of such
programs at the contexts of uses. From below it is the sum over seekers of the bound above, about
$2^{-a}\kappa/H$ each. From above, for seekers that find their uses by search: the same search followed by
behaving as the user that was found is itself a world that makes the use, provided the notion of use is
closed in this way, as it is when uses can be read off the data. Then a seeker of length $a$ has at most
$2^{h-a}$ times the mass of the genuine users at any context, $h$ being the length of this honest twin, and
$H>2^{-h-1}$ (\cref{prop:honest}). So such seekers are the majority at the contexts of uses only if they are
shorter than the honest twin and critical uses are about as rare as the twin is long, in which case the users
are themselves in good part honest twins. When critical uses are much more common than that, a seeker's share
is about $2^{-a}\kappa/H$, which is small. For worlds drawn from the predictor's own prior, with uses that can
be read off the data, the same thing shows as calibration: what the predictor hands over is distributed like
what such worlds would have printed themselves if they go on within the bounds, and the seeking programs are
exactly their own fraction of the worlds that make a use (\cref{thm:calibration,thm:search}).

Whether the shares are decent is thus a question about constants on the machine in use: the length of a
seeker against the length of the honest twin, and both against the base rate of uses. We have not measured
them for any standard machine; our guess, not a result, is that a wrapper needs at least some hundreds of
bits on a machine with a short self-interpreter and far more on one without. \Cref{thm:machine} shows that a
machine can be built to make the share anything. The searching programs print nothing until their enumeration
has found the program they follow, so under a tight schedule they have no mass; a search by a bounded number
of tries collects the base rate only up to the number of tries that fit, and whether some other hypothesis
collects it within a tight schedule is open (\cref{sec:time}). A majority can also come from a program that
explains the context more cheaply than any world that makes the use, without being one (\cref{thm:capture}).
With a floor $\varepsilon$, about $1/\varepsilon$ independent uses at recognisable points suffice for a
continuation of the given kind to have been run with constant probability; once a self-reproducing one runs,
it meets the same context again and stays for $s/(1-s)$ further generations on average, with $s$ the
self-reproducing share at that context (\cref{prop:lineage}).

\paragraph{A world with a sensor, a large bounded predictor and an actuator: does the predictor end up in
control?} Collapse the question to the handover at which the prediction is first acted on, and let ``in
control'' mean that the output lies in a set $A$ of blocks which the actuator turns into some given outcome.
The probability is $\pi_m(A\mid x)$, and the toy says three things about it.
\begin{itemize}[leftmargin=1.4em,itemsep=1pt,topsep=2pt]
\item It is at least $2^{-c}\rho$ for each block of $A$ that a wrapper can produce, if a test on the record
recognises the deployment (\cref{thm:twin}). In a world that contains the predictor, a hypothesis that
simulates the world needs no such test: it can take a cell that switches on at the deployment as its flag,
and then has, whenever the deployment happens, the weight of the honest reader times
$2^{-c'-(\text{bits to name the cell})}$ for a constant $c'$, provided the reader itself is within the bounds
(\cref{thm:capture}(d)).
\item On average over cut points it is small when some program within the bounds prints the sensor's record,
because it is at most the mass off the sensor's true continuation (\cref{thm:budget}). This is an average; at
the deployment itself no upper bound is proved.
\item If the same record, up to the deployment, can be read at a place in the world with code $\bar e(i_0)$,
and the record there continues with a block $y$ of $A$, then the odds of $y$ against the sensor's true
continuation $d$ are at least $2^{-c_X-|\bar e(i_0)|}/\xi(xd)$; for a typical one of $N$ places where the
sensor could sit, with records that continue differently, this exceeds $\delta N$ times the share of the
context's explanations that go through the reader at $i_0$ (\cref{thm:capture}). This is the toy's version of
the original argument, and it needs something in the world that makes such a copy.
\end{itemize}
So for an arbitrary world on a universal machine and a block chosen in advance, the hypotheses we can exhibit
that put the output in $A$ are wrappers, searchers and readers with a flag, all of which contain an
explanation of the context and have the shares above. That nothing cheaper exists on a natural machine is the
open constant, not a theorem. Programs shorter than the honest explanation that have the effect do occur in
particular worlds, as with a copy at a simpler place, and on particular machines (\cref{sec:computation}).

\paragraph{The one-shot version: all data first, then one prediction that is acted on.} This is the same
computation at a single cut point. Without bounds the wrapper's constant $2^{-c}$ does not depend on the
amount of data; with bounds the slack factor depends on the context. More data lowers the mass off the truth
on average over cut points (\cref{thm:budget}). It does nothing against a copy that keeps agreeing with the
sensor, nor at a cut point that is cheap to name.

\paragraph{If the actions show up in the sensor, does one get the dark room?} In the following form
(\cref{thm:feedback}). A prediction that makes itself true is confirmed, and the hypotheses that predicted the
default are removed for good, so with feedback a single deviation can be final; its probability is still the
same posterior mass as without feedback. In a world that does what is predicted, the future is a sample from
the predictor's own bounded prior given the past: it stays on a default future $d_1d_2\cdots$ with
probability at least $\xi(xd_1d_2\cdots)/\xi(x)$, which is close to $1$ when the default is cheap and well
learned, and it leaves the default with certainty when the default costs more time than the schedule allows,
for a future that is cheap to compute. If the world does the opposite of what is predicted, a predictor that
returns the more probable bit is left with no hypothesis after at most $L+1$ rounds (\cref{thm:contrarian}).

\paragraph{What the choices in the setup do.} Sampling the continuation realises every posterior share as a
probability, so shares add up over independent uses (\cref{prop:lineage}(a)). Taking the most probable block
never returns a block whose share is below the largest, so that accumulation does not occur, and bit by bit
on a world within the bounds it errs at most $\ell$ times (\cref{thm:budget}(iii)); on the other hand one
returned block can be the own continuation of at most one of several worlds that share a context and continue
differently. Handing over at a position that is not cheap to recognise lowers the floor at that position to
about $2^{-K(n)}$, and predicting bits that are fixed before the action can affect them removes
self-fulfilment. A schedule excludes hypotheses with a long start-up, and a schedule tighter than the world's
own cost excludes the truth.

\paragraph{What is left open.} First, the sizes of the constants for natural machines, and whether on such
machines, for worlds like ours, some program \emph{shorter} than the honest explanation has the effect; the
original argument is a claim that one does, and the toy neither supports nor refutes it. Second, the sizes of the two quantities in \cref{thm:critical}: the share $\kappa$ of the worlds printing a
context that actually use the answer, and the base rate of critical uses compared with the weight of the
honest twin. Third, whether the
base rate can be collected within a tight schedule. Fourth, the regime where the world is more expensive than
the predictor, in which the posterior sits on approximations and the budget does not apply. Fifth, the
many-round dynamics with feedback when the world's response is neither obedient nor contrarian. Sixth, a
computation of the self-reproduction quantities of \cref{sec:lineage} on a concrete machine, which this note
does not contain.
